\subsection{The arithmetic free-inertia comparison and its correction words}

\begin{statement}{E.39}{Proved}{Riemann existence gives the exact kernel, with unavoidable Frobenius corrections}
Put $K=\mathbb Q(\zeta)$, $S=\{(1-\zeta)\}$,
$\mathcal G=G_K^{(3)}$, and
$P=G_{K,S}^{(3)}=\langle\tau,\tau'\rangle$ as in E.37.
Let $L=\ker(\mathcal G\to P)$. There is an actual
continuous free-product isomorphism
\[
 \underset{w\nmid3\text{ in }K_S(3)}{\mathop{\ast}} I_w
       \ \xrightarrow{\ \sim\ } L,
 \qquad I_w\simeq\mathbb Z_3(1),
 \tag{E39.1}
\]
induced by the inclusions of the chosen inertia groups.
The indexing and free product have the continuous-bundle
topology; finite quotients retain only finitely many base
primes and locally constant data on their prolongations.
Choose lifts of the two Frobenii over seven, related by
complex conjugation. They define a conjugation-equivariant
splitting $s:P\to\mathcal G$, and hence
$\mathcal G=L\rtimes s(P)$.

For an actual local Frobenius $F_{\mathfrak q}$ put
\[
 \sigma_{\mathfrak q}=\operatorname{pr}(F_{\mathfrak q}),
 \qquad h_{\mathfrak q}
       =F_{\mathfrak q}s(\sigma_{\mathfrak q})^{-1}\in L.
 \tag{E39.2}
\]
No splitting makes every $h_{\mathfrak q}$ belong to
the corresponding local inertia group. For the specified
splitting, at $\mathfrak q=(-2+3\zeta)$ above nineteen,
the actual cyclic cubic quotient already detects
\[
 \kappa_{2,3}(h_{\mathfrak q})=2\in\mathbb F_3,
 \tag{E39.3}
\]
where $\kappa_{2,3}$ is the Kummer character of $2$.
This character is zero on $I_{\mathfrak q}$.

The precise section equations supplied by (E39.1) are
as follows. Let $W$ be $N$, $N/N''$, or one of the
Galois-stable marked quotients in E.25, with its actual
based action. A section over $K$ is specified by a
continuous homomorphism $\lambda:L\to W$ and a cocycle
$u:P\to W$ satisfying
\[
 \lambda(s(g)\ell s(g)^{-1})
       =u(g)\,g(\lambda(\ell))\,u(g)^{-1}.
 \tag{E39.4}
\]
Here $u$ is freely specified by its two values on
$\tau,\tau'$, and $\lambda$ is specified by its continuous
family of values on the factors in (E39.1). The section
cocycle is reconstructed by the exact formula
\[
 a(\ell s(g))=\lambda(\ell)u(g).
 \tag{E39.5}
\]
For tame inertia $t_{\mathfrak q}$, writing
$n_{\mathfrak q}=\lambda(t_{\mathfrak q})$ and
$q=N\mathfrak q$, (E39.4) gives the necessary local equation
\[
 [\lambda(h_{\mathfrak q})u(\sigma_{\mathfrak q})]
       \sigma_{\mathfrak q}(n_{\mathfrak q})
 [\lambda(h_{\mathfrak q})u(\sigma_{\mathfrak q})]^{-1}
       =n_{\mathfrak q}^{q}.
 \tag{E39.6}
\]
The correction word is part of this formula, including
for metabelian $W$.

\emph{Manuscript: \S51.16, folios 733v--736v, scans 172--178;
Part E(c),(d). These are comparisons with the actual arithmetic
Galois group and intrinsic geometric quotients.}
\end{statement}
\begin{proof}
The precise primary result used here is Wingberg,
\cite[Theorem 2.4, pp.~13--14]{E_WingbergRiemann}; its
forward implication explicitly includes infinite sets of
primes. The $K(\pi,1)$ hypothesis holds for
$\operatorname{Spec}\mathcal O_K[1/3]$, as stated in the
remark after Proposition 2.1 of that paper. It also follows
from the finite-coefficient calculation in E.34: its
degree-two cohomology vanishes, and induction along finite
three-covers preserves that vanishing. The universal
three-cover kills degree-one cohomology. The local
nontriviality hypothesis holds at every $\mathfrak q\nmid3$
because the cyclotomic $\mathbb Z_3$-extension is inside
$K_S(3)$ and its arithmetic Frobenius is $q\ne1$ in
$1+3\mathbb Z_3$. The theorem therefore gives (E39.1)
for the set of all finite primes. At these residue
characteristics, pro-three inertia is $\mathbb Z_3(1)$.

The quotient $P$ is the free rank-two group already
proved in E.34 and E.37. Select an actual Frobenius lift
above $(1+3\zeta)$ in $\mathcal G$, and use its complex
conjugate above the other prime over seven. The free
property produces $s$ with these values. Its composite
with $\mathcal G\to P$ fixes both free generators and
is the identity. Conjugation interchanges the two chosen
lifts, so $s$ also commutes with conjugation.

For completeness, the failure of simultaneous local
normalization holds for every splitting, not only this
one. Let $\chi=\kappa_{2,3}:\mathcal G\to\mathbb F_3$,
and set $\beta=\chi\circ s:P\to\mathbb F_3$.
The character $\chi$ ramifies at the prime $(2)$:
the local polynomial $X^3-2$ is Eisenstein, and $K$
contains $\mu_3$. The character $\beta$ is unramified
away from three. Were all $h_{\mathfrak q}$ local-inertia
elements, for every prime away from six we would have
$\chi(F_{\mathfrak q})=\beta(\sigma_{\mathfrak q})$.
Apply the abelian Chebotarev theorem to the finite
extension cut out by these two characters
\cite[V, Theorem 3.23]{E_MilneCFT}.
It makes the characters equal on all of $\mathcal G$,
contradicting their different inertia values at two.
This proves the asserted failure by actual character maps.

There is also the finite witness (E39.3), requiring no
unspecified prime. At the primes over seven, the images
of $\zeta$ in the residue fields are two and four.
For a unit $b$, the exponent of $\zeta$ in
$b^{(q-1)/3}$ is its cubic Kummer Frobenius value. These
exponents give
\[
 \begin{array}{c|ccc}
    &\kappa_{2,3}&\kappa_{\zeta,3}&\kappa_{\lambda,3}\\\hline
 \zeta\equiv2\pmod7&2&2&0\\
 \zeta\equiv4\pmod7&1&2&2.
 \end{array}
 \tag{E39.7}
\]
The last two characters are the two unramified Kummer
generators, and evaluation on the free generators shows
\[
   \kappa_{2,3}\circ s=\kappa_{\zeta\lambda,3}
          \quad\text{as characters of }P.
 \tag{E39.8}
\]
At $\mathfrak q=(-2+3\zeta)$ the residue of $\zeta$ is
seven, that of $\lambda$ is thirteen, and
$2^6=7$, $(\zeta\lambda)^6=11=7^2$ modulo nineteen.
Thus (E39.2) has character value $1-2=2$ modulo three.
The program \texttt{e\_riemann\_correction.py} verifies
all residues and the two-character calculation.

The based action on $W$ factors through $P$ over $K$,
by E.34. A section cocycle factors through
$\mathcal G$: the homomorphism into $W$ semidirect
the arithmetic image has pro-three target. Its
restriction to $L$ is a homomorphism because $L$ acts
trivially on $W$. Restriction to $s(P)$ is the cocycle
$u$. Applying the cocycle identity to $s(g)\ell s(g)^{-1}$
gives (E39.4). Conversely this identity verifies the
cocycle law for (E39.5), using
$\mathcal G=L\rtimes s(P)$, so the displayed description
is a bijection of actual cocycles. The free-product
universal property gives exactly the continuity
condition stated for $\lambda$; arbitrary discontinuous
families are excluded.

For sections over $\mathbb Q$, first conjugate so that
the value at complex conjugation $c$ is one. This is
possible because the complement of the pro-three kernel
in its extension by $\langle c\rangle$ is conjugate
to the standard complement, by the finite
Schur--Zassenhaus theorem and inverse limits.
Then add the two exact equations
\[
  u(\tau')=c(u(\tau)),\qquad
  \lambda(c\ell c)=c(\lambda(\ell)).
 \tag{E39.9}
\]
The remaining geometric conjugations are $w\in W^c$;
they send $(\lambda,u)$ to
$(w\lambda w^{-1},\,g\mapsto w u(g)g(w)^{-1})$.
This specifies the rational descent and its conjugacy
quotient rather than changing the arithmetic base field.

Finally $F_{\mathfrak q}=h_{\mathfrak q}s(\sigma_{\mathfrak q})$
and the actual local relation is
$F_{\mathfrak q}t_{\mathfrak q}F_{\mathfrak q}^{-1}
=t_{\mathfrak q}^q$. Apply (E39.5) to obtain (E39.6).
In a metabelian target the value $\lambda(h_{\mathfrak q})$
still uses the image of that word in $L/L''$; (E39.3)
shows that even its first abelian coordinate cannot
be set to zero in general. Equations (E39.4)--(E39.9)
therefore record the precise further arithmetic input
needed for this approach. They do not replace the
evaluation of those words, or of their conjugation
action, by the two unramified Frobenius matrices.
\end{proof}
